\documentclass[11pt,a4paper]{article}

\usepackage[utf8]{inputenc}
\usepackage[T1]{fontenc}
\usepackage[spanish,es-tabla]{babel}
\usepackage{lmodern}
\usepackage{geometry}
\usepackage{array,tabularx,booktabs}
\usepackage[table]{xcolor}
\usepackage{enumitem}
\usepackage{titlesec}
\usepackage{microtype}
\usepackage[hidelinks]{hyperref}

\geometry{left=1.28cm,right=1.28cm,top=1.08cm,bottom=1.08cm}
\setlength{\parindent}{0pt}
\setlength{\parskip}{0.09em}
\setlist[itemize]{leftmargin=1.0cm,itemsep=0.02em,topsep=0.05em}
\setlist[enumerate]{leftmargin=0.92cm,itemsep=0.01em,topsep=0.04em}
\titleformat{\section}{\normalsize\bfseries}{\thesection.}{0.4em}{}
\titlespacing*{\section}{0pt}{0.38em}{0.08em}
\renewcommand{\baselinestretch}{0.945}
\pagestyle{empty}

\newcommand{\universidad}{Universidad Cat\'olica ``Nuestra Se\~nora de la Asunci\'on''}
\newcommand{\facultad}{Facultad de Ciencias y Tecnolog\'ias}
\newcommand{\departamento}{Departamento de Ingenier\'ia Electr\'onica e Inform\'atica}
\newcommand{\titulo}{Dise\~no e implementaci\'on de un sistema electr\'onico multicanal del tipo IoT para caracterizaci\'on geot\'ecnica del suelo mediante ondas s\'ismicas}
\newcommand{\alumno}{El\'ias David \'Alvarez Mart\'inez}
\newcommand{\carrera}{Ingenier\'ia Electr\'onica}
\newcommand{\tutor}{Enrique A. Vargas Cabral, Ph.D.}
\newcommand{\matricula}{Y24137}

\begin{document}

\begin{center}
    {\large \universidad}\par
    {\facultad}\par
    {\departamento}\par
    \vspace{0.28cm}
    {\Large\bfseries \titulo}\par
    \vspace{0.14cm}
    {\bfseries Propuesta de Proyecto Final de Carrera}\par
    Junio 2026
\end{center}
\vspace{0.07cm}
\begin{tabularx}{\textwidth}{@{}lX lX@{}}
\textbf{Alumno:} & \alumno & \textbf{Carrera:} & \carrera \\
\textbf{Tutor:} & \tutor & \textbf{Matr\'icula:} & \matricula \\
\end{tabularx}

\section{Objetivo General}
Dise\~nar, implementar y evaluar experimentalmente un sistema electr\'onico multicanal para la generaci\'on y captaci\'on de ondas s\'ismicas, orientado a la caracterizaci\'on geot\'ecnica del suelo en profundidades m\'aximas de hasta 50 metros. El sistema comprender\'a la selecci\'on e implementaci\'on del generador de ondas s\'ismicas, el dise\~no de ge\'ofonos aut\'onomos del tipo IoT (Internet of Things), su organizaci\'on en configuraciones matriciales de medici\'on y el procesamiento digital de las se\~nales adquiridas. A partir de dichas mediciones se buscar\'a estimar perfiles de velocidad de propagaci\'on y par\'ametros el\'asticos asociados a las distintas subcapas del suelo, con el fin de disponer de una herramienta experimental complementaria para estudios geot\'ecnicos.

\section{Justificaci\'on y metodolog\'ia}
La caracterizaci\'on geot\'ecnica del suelo es una etapa fundamental para el dise\~no de cimentaciones, obras civiles, estructuras sometidas a cargas din\'amicas y proyectos que requieran conocer el comportamiento mec\'anico del terreno \cite{Kramer1996,Mitchell2005}. En la pr\'actica profesional se utilizan ensayos tradicionales como el Standard Penetration Test (SPT) y el Cone Penetration Test (CPT), los cuales entregan informaci\'on valiosa, pero pueden implicar intervenci\'on directa del terreno, tiempos de ejecuci\'on considerables, costos operativos elevados y dificultades log\'isticas en zonas de acceso complejo \cite{ASTM}. Estas limitaciones motivan el estudio de m\'etodos complementarios, de menor impacto sobre el suelo y con potencial para realizar mediciones m\'as r\'apidas en campo.

Entre las t\'ecnicas geof\'isicas aplicables se destacan la Tomograf\'ia de Refracci\'on S\'ismica (SRT) y el An\'alisis Multicanal de Ondas Superficiales (MASW). La SRT permite estimar perfiles de velocidad a partir de los tiempos de llegada de ondas refractadas, mientras que el MASW utiliza la dispersi\'on de ondas superficiales, principalmente ondas Rayleigh, para obtener informaci\'on relacionada con la velocidad de ondas de corte \cite{Park1999,Xia1999,Foti2014}. La literatura reciente sobre MASW recomienda la adquisici\'on con arreglos multicanal y m\'ultiples configuraciones fuente-receptor para mejorar la resoluci\'on, extender el rango de frecuencias y evaluar la incertidumbre de las curvas de dispersi\'on \cite{Foti2018,OlafsdottirNGM2024}. Por ello, el desarrollo de un sistema propio de adquisici\'on multicanal permite integrar instrumentaci\'on electr\'onica, procesamiento digital y validaci\'on experimental en un mismo prototipo.

Desde el punto de vista electr\'onico, el ge\'ofono y su acondicionamiento son elementos cr\'iticos, ya que la frecuencia natural, el amortiguamiento, el ruido del circuito y la sincronizaci\'on temporal influyen directamente en la calidad de las se\~nales registradas \cite{Ma2023,Olafsdottir2018}. La implementaci\'on de nodos aut\'onomos del tipo IoT busca reducir el cableado, simplificar el despliegue en campo y permitir configuraciones matriciales flexibles, manteniendo el registro coordinado de eventos s\'ismicos de corta duraci\'on. En este contexto, herramientas abiertas como MASWavesPy constituyen una referencia \'util para procesar datos MASW, comparar resultados e implementar rutinas de inversi\'on de perfiles de velocidad \cite{Olafsdottir2020,OlafsdottirNGM2024}.

La metodolog\'ia partir\'a del estudio del estado del arte sobre propagaci\'on de ondas s\'ismicas, ge\'ofonos, sistemas multicanal y adquisici\'on inal\'ambrica. Posteriormente se seleccionar\'an los componentes principales del sistema, incluyendo el mecanismo de excitaci\'on s\'ismica, sensores, acondicionamiento anal\'ogico, conversi\'on anal\'ogico-digital, unidad de procesamiento, alimentaci\'on y comunicaci\'on. Luego se dise\~nar\'an nodos ge\'ofonos aut\'onomos con sincronizaci\'on temporal suficiente para registrar eventos s\'ismicos de corta duraci\'on. Finalmente, se integrar\'a el prototipo, se ejecutar\'an ensayos experimentales y se aplicar\'an algoritmos de filtrado, detecci\'on de tiempos de llegada, an\'alisis de dispersi\'on e inversi\'on de velocidades. Los resultados ser\'an contrastados con datos obtenidos mediante ensayos SPT o informaci\'on geot\'ecnica disponible.

\section{Etapas de ejecuci\'on del trabajo}
\begin{enumerate}
    \item Estudio de los principios de propagaci\'on de ondas s\'ismicas en suelos y revisi\'on del estado del arte sobre SRT, MASW, ge\'ofonos, adquisici\'on multicanal y sistemas IoT aplicados a instrumentaci\'on de campo.
    \item Selecci\'on del generador de ondas s\'ismicas, de los ge\'ofonos y de los componentes electr\'onicos principales para la captaci\'on, acondicionamiento, digitalizaci\'on y transmisi\'on de las se\~nales.
    \item Dise\~no de los circuitos de acondicionamiento, alimentaci\'on, sincronizaci\'on y registro de datos para los nodos ge\'ofonos aut\'onomos.
    \item Implementaci\'on del prototipo electr\'onico multicanal y definici\'on de las topolog\'ias de matrices de ge\'ofonos a emplear durante las pruebas de campo.
    \item Desarrollo de los algoritmos de procesamiento digital de se\~nales, incluyendo filtrado, detecci\'on de llegadas, an\'alisis de dispersi\'on e inversi\'on para estimaci\'on de perfiles de velocidad.
    \item Integraci\'on del sistema, realizaci\'on de ensayos experimentales, comparaci\'on de resultados con el ensayo SPT, an\'alisis de desempe\~no y redacci\'on de la memoria final del Proyecto Final de Carrera.
\end{enumerate}

\section{Cronograma de actividades}
\begin{center}
\scriptsize
\renewcommand{\arraystretch}{1.12}
\setlength{\tabcolsep}{2.6pt}
\begin{tabularx}{0.90\textwidth}{|>{\raggedright\arraybackslash}X|*{6}{>{\centering\arraybackslash}p{0.70cm}|}}
\hline
\rowcolor{gray!20}
 & \multicolumn{6}{c|}{\textbf{Per\'iodos (en meses)}} \\ \cline{2-7}
\rowcolor{gray!20}
\textbf{Actividad} & \textbf{M1} & \textbf{M2} & \textbf{M3} & \textbf{M4} & \textbf{M5} & \textbf{M6} \\ \hline
Estudio del estado del arte. & \cellcolor{gray!35} & & & & & \\ \hline
Definici\'on y dise\~no de la arquitectura y los circuitos el\'ectricos. & & \cellcolor{gray!35} & \cellcolor{gray!35} & & & \\ \hline
Dise\~no, simulaci\'on e implementaci\'on de la estructura f\'isica. & & & \cellcolor{gray!35} & \cellcolor{gray!35} & & \\ \hline
Integraci\'on del prototipo. & & & & & \cellcolor{gray!35} & \\ \hline
Montaje y evaluaci\'on del prototipo. & & & & & & \cellcolor{gray!35} \\ \hline
Redacci\'on de la memoria del TFG. & \cellcolor{gray!35} & \cellcolor{gray!35} & \cellcolor{gray!35} & \cellcolor{gray!35} & \cellcolor{gray!35} & \cellcolor{gray!35} \\ \hline
\end{tabularx}
\par\vspace{0.25em}
{\small Figura 1: Cronograma de actividades.}
\end{center}

\vspace{-0.45em}
\section{Referencias}
\renewcommand{\refname}{}
\vspace{-1.7em}
\begin{thebibliography}{99}\footnotesize
\setlength{\itemsep}{0.02em}
\bibitem{Kramer1996} S. L. Kramer, \textit{Geotechnical Earthquake Engineering}. Prentice Hall, 1996.
\bibitem{Foti2014} S. Foti, C. G. Lai, G. J. Rix y C. Strobbia, \textit{Surface Wave Methods for Near-Surface Site Characterization}. CRC Press, 2014.
\bibitem{Park1999} C. B. Park, R. D. Miller y J. Xia, ``Multichannel analysis of surface waves'', \textit{Geophysics}, vol. 64, no. 3, pp. 800--808, 1999.
\bibitem{Xia1999} J. Xia, R. D. Miller y C. B. Park, ``Estimation of near-surface shear-wave velocity by inversion of Rayleigh waves'', \textit{Geophysics}, vol. 64, no. 3, pp. 691--700, 1999.
\bibitem{Mitchell2005} J. K. Mitchell y K. Soga, \textit{Fundamentals of Soil Behavior}. John Wiley \& Sons, 2005.
\bibitem{ASTM} ASTM International, \textit{Standard Test Method for Standard Penetration Test (SPT) and Split-Barrel Sampling of Soils}, ASTM D1586/D1586M.
\bibitem{Ma2023} K. Ma, J. Wu, Y. Ma, B. Xu, S. Qi y X. Jiang, ``An Effective Method for Improving Low-Frequency Response of Geophone'', \textit{Sensors}, vol. 23, no. 6, art. 3082, 2023. doi:10.3390/s23063082. Disponible: \url{https://pmc.ncbi.nlm.nih.gov/articles/PMC10059280/}.
\bibitem{OlafsdottirNGM2024} E. A. Olafsdottir, B. Bessason y S. Erlingsson, ``MASWavesPy: A Python Package for Analysis of MASW Data'', \textit{19th Nordic Geotechnical Meeting}, G\"oteborg, 2024. Disponible: \url{https://ngm2024.se/Papers/ngm2024_maswavespy_final.pdf}.
\bibitem{Olafsdottir2018} E. A. Olafsdottir, S. Erlingsson y B. Bessason, ``Tool for analysis of multichannel analysis of surface waves field data and evaluation of shear wave velocity profiles of soils'', \textit{Canadian Geotechnical Journal}, vol. 55, no. 2, pp. 217--233, 2018. doi:10.1139/cgj-2016-0302.
\bibitem{Olafsdottir2020} E. A. Olafsdottir, S. Erlingsson y B. Bessason, ``Open-Source MASW Inversion Tool Aimed at Shear Wave Velocity Profiling for Soil Site Explorations'', \textit{Geosciences}, vol. 10, no. 8, 2020. doi:10.3390/geosciences10080322.
\bibitem{Foti2018} S. Foti et al., ``Guidelines for the good practice of surface wave analysis: a product of the InterPACIFIC project'', \textit{Bulletin of Earthquake Engineering}, vol. 16, no. 6, pp. 2367--2420, 2018. doi:10.1007/s10518-017-0206-7.
\bibitem{Vantassel2022} J. P. Vantassel y B. R. Cox, ``SWprocess: a workflow for developing robust estimates of surface wave dispersion uncertainty'', \textit{Journal of Seismology}, vol. 26, pp. 731--756, 2022. doi:10.1007/s10950-021-10035-y.
\bibitem{Wathelet2020} M. Wathelet et al., ``Geopsy: a user-friendly open-source tool set for ambient vibration processing'', \textit{Seismological Research Letters}, vol. 91, no. 3, pp. 1878--1889, 2020. doi:10.1785/0220190360.
\bibitem{Tokimatsu1992} K. Tokimatsu, S. Tamura y H. Kojima, ``Effects of multiple modes on Rayleigh wave dispersion characteristics'', \textit{Journal of Geotechnical Engineering}, vol. 118, no. 10, pp. 1529--1543, 1992. doi:10.1061/(ASCE)0733-9410(1992)118:10(1529).

\bibitem{Wood2012} C. M. Wood y B. R. Cox, ``A comparison of MASW dispersion uncertainty and bias for impact and harmonic sources'', \textit{Proc. of GeoCongress 2012}, pp. 2756--2765, 2012. doi:10.1061/9780784412121.282.
\bibitem{Buchen1996} P. W. Buchen y R. Ben-Hador, ``Free-mode surface-wave computations'', \textit{Geophysical Journal International}, vol. 124, no. 3, pp. 869--887, 1996. doi:10.1111/j.1365-246X.1996.tb05642.x.
\bibitem{Long2007} M. Long y S. Donohue, ``In situ shear wave velocity from multichannel analysis of surface waves (MASW) tests at eight Norwegian research sites'', \textit{Canadian Geotechnical Journal}, vol. 44, no. 5, pp. 533--544, 2007. doi:10.1139/t07-013.
\bibitem{Blaker2019} \O. Blaker et al., ``Halden research site: geotechnical characterization of a post glacial silt'', \textit{AIMS Geosciences}, vol. 5, no. 2, pp. 184--234, 2019. doi:10.3934/geosci.2019.2.184.
\bibitem{Quinteros2019} S. Quinteros et al., ``\O{}ysand research site: Geotechnical characterisation of deltaic sandy-silty soils'', \textit{AIMS Geosciences}, vol. 5, no. 4, pp. 750--783, 2019. doi:10.3934/geosci.2019.4.750.
\end{thebibliography}

\vspace{0.15em}
\begin{flushright}
\rule{6.0cm}{0.4pt}\\[-0.1em]
{\small Vo.Bo. \tutor}
\end{flushright}

\end{document}
